% FILE: 03_architecture_and_primitives.tex
% Project: otel-weaver-exp-002-diff-to-residual
% Thesis Role: weaver-diff-to-pi-residual-bridge-experiment

\section{Architecture and Primitives}
\label{sec:otel-weaver-exp-002-diff-to-residual-arch}

\subsection{Core Objects}
\label{sec:otel-weaver-exp-002-diff-to-residual-arch-objects}

The experiment defines five core objects in \texttt{src/bridge.rs}:

\paragraph{TelemetryFeedstockS2}
A \texttt{serde}-deserializable Rust struct representing telemetry conforming to
schema version S2 of the process intelligence semantic conventions. It carries four
fields:
\begin{itemize}
    \item \texttt{instance\_id}: mapped from \texttt{process.pi.instance\_id}
    \item \texttt{transition\_name}: mapped from \texttt{process.pi.transition.name}
          (the renamed field)
    \item \texttt{token\_color}: mapped from \texttt{process.pi.token.color} as
          \texttt{Option<String>} (the added field, supporting Colored Petri Nets)
    \item \texttt{witness\_hash}: mapped from \texttt{process.pi.witness.hash}
\end{itemize}

\paragraph{TelemetryFeedstockS1}
A \texttt{serde}-deserializable Rust struct representing telemetry conforming to
schema version S1. It carries three fields:
\begin{itemize}
    \item \texttt{instance\_id}: mapped from \texttt{process.pi.instance\_id}
    \item \texttt{activity\_name}: mapped from \texttt{process.pi.activity.name}
          (the original field name, expected by the conformance court)
    \item \texttt{witness\_hash}: mapped from \texttt{process.pi.witness.hash}
\end{itemize}
The absence of \texttt{token\_color} in S1 is intentional: the S1 schema predates
Colored Petri Net token tracking. What happens to the \texttt{token\_color} value
during translation is an open question documented in
Section~\ref{sec:otel-weaver-exp-002-diff-to-residual-open}.

\paragraph{BridgeRx}
The central translation struct. It holds a
\texttt{rename\_map: HashMap<String, String>} that registers known field renames from
S2 to S1. At construction (\texttt{BridgeRx::new()}), the map is populated with the
single known rename:
\begin{verbatim}
"process.pi.transition.name" -> "process.pi.activity.name"
\end{verbatim}
The \texttt{translate(\&self, input: TelemetryFeedstockS2) -> TelemetryFeedstockS1}
method applies the rename, preserves the witness hash, and drops \texttt{token\_color}.

\paragraph{PI Conformance Residual $\mathcal{R}_{\Delta}$}
The mathematical proof target. Defined in \texttt{README.md} as:
\[
    \mathcal{R}_{\Delta} = |f_{S_1}(\mathcal{L}_2, M) - f_{S_2}(\mathcal{L}_2, M)|
\]
After \texttt{BridgeRx} translation, $\mathcal{R}_{\Delta} = 0$ because the
translated S1 record is evaluated by the S1 parser, eliminating the fitness
discrepancy between parser versions.

\paragraph{Weaver Diff $\Delta_W(S_1, S_2)$}
A set of \texttt{(path, op, value)} tuples where
$op \in \{\text{Add}, \text{Remove}, \text{Rename}, \text{TypeChange}\}$.
The experiment's worked example in \texttt{README.md} shows a two-tuple diff:
a \texttt{Rename} on \texttt{process.pi.activity.name} and an \texttt{Add} on
\texttt{process.pi.token.color}.

\subsection{Admissible Transitions}
\label{sec:otel-weaver-exp-002-diff-to-residual-arch-transitions}

The \texttt{BridgeRx} architecture admits only one lawful transition path from
S2 feedstock to the conformance boundary: through \texttt{translate()}. The Rust
type system enforces this: the conformance court's validators accept
\texttt{TelemetryFeedstockS1}, not \texttt{TelemetryFeedstockS2}. There is no
implicit coercion between the two types; the only path from one to the other is
the explicit \texttt{bridge.translate(s2\_record)} call.

This is an architectural enforcement of the categorical law
$\Delta_W \neq \delta_P$ at the type level: S2 records simply cannot be evaluated
for process conformance without passing through the bridge. The Rust compiler
refuses to compile code that attempts to pass an S2 struct where an S1 struct is
expected, making the category separation a build-time guarantee rather than a
runtime convention.

The upstream \texttt{Admission<T,W>} and \texttt{Refusal} type-law patterns
(documented in exp-003 and referenced in \texttt{GGEN\_OTEL\_WEAVER\_PI\_ALIVE\_001.md})
represent the next admissible transition after bridging: a translated S1 record
may enter the \texttt{Admission} gate; a record that fails the gate emits a
\texttt{Refusal} with a cryptographic BLAKE3 digest of the failure context.

\subsection{Boundaries}
\label{sec:otel-weaver-exp-002-diff-to-residual-arch-boundaries}

Two boundary rules are formally defined in
\texttt{mappings/registry-diff-to-pi-drift-map.yaml}:

\paragraph{drift\_isolation\_rule\_1}
Condition: a telemetry field is missing because of a schema diff (e.g., attribute
deprecation). Then: do not report process drift; route to the Weaver Diff Log at
\texttt{checkpoints/weaver\_diff\_registry.json}.

\paragraph{drift\_isolation\_rule\_2}
Condition: a telemetry field is missing but is present in the current schema version.
Then: report Process Drift (Activity Skip / Data Under-provisioning) to the process
drift court at \texttt{checkpoints/process\_drift\_court.json}.

These two rules together define the decision tree that the ingestion layer must
execute before any conformance algorithm fires. \texttt{BridgeRx} is the
implementation of \texttt{drift\_isolation\_rule\_1} for the \texttt{Rename} case:
it ensures that a field missing from S1 terms but present in S2 terms is not routed
to the drift court.

The \texttt{otel-to-pi-witness-map.yaml} mapping file defines the boundary rules for
translating OTel service identity to PI witness IDs, ensuring that the
\texttt{pi.feedstock.witness.id} field is correctly populated as telemetry crosses
the ingestion boundary. The mapping also enforces the conformance isolation
prevention rule: ``Never map feedstock directly to court consequences.''

\subsection{Failure Modes Refused}
\label{sec:otel-weaver-exp-002-diff-to-residual-arch-failures}

The architecture of EXP-002 is designed to refuse three specific failure modes:

\begin{enumerate}
    \item \textbf{Measurement theater}: Schema naming updates flagged as compliance
          violations. Refused by \texttt{BridgeRx} translating S2 field names to S1
          expectations before the conformance algorithm evaluates fitness.

    \item \textbf{Categorical collapse}: The false inference that
          $\Delta_W(S_1, S_2) = \delta_P(\mathcal{L}_1, \mathcal{L}_2)$. Refused
          by the type system (S2 and S1 are distinct types) and by the doctrine
          files anchored at \texttt{REGISTRY\_DIFF\_NOT\_DRIFT\_001} and
          \texttt{OTEL\_WEAVER\_FEEDSTOCK\_001}.

    \item \textbf{Witness hash corruption}: The BLAKE3 hash emitted by
          instrumentation is altered during schema translation. Refused by the
          \texttt{translate()} implementation's direct field assignment
          \texttt{witness\_hash: input.witness\_hash} with no transformation, and
          asserted by \texttt{test\_bridge\_translation\_equivalence}.
\end{enumerate}

The \texttt{doctrine/registry-diff-is-not-process-drift.md} file also documents a
fourth failure mode this architecture prevents: the \textbf{false negative}, where
actual process drift (e.g., bypass of approval gates) goes undetected because the
telemetry payload is structurally clean under the current Weaver schema. By keeping
the conformance court independent of schema validation, genuine process deviations
remain visible even when the telemetry surface is structurally compliant.
